%---------- Inleiding ---------------------------------------------------------

% TODO: Is dit voorstel gebaseerd op een paper van Research Methods die je
% vorig jaar hebt ingediend? Heb je daarbij eventueel samengewerkt met een
% andere student?
% Zo ja, haal dan de tekst hieronder uit commentaar en pas aan.

%\paragraph{Opmerking}

% Dit voorstel is gebaseerd op het onderzoeksvoorstel dat werd geschreven in het
% kader van het vak Research Methods dat ik (vorig/dit) academiejaar heb
% uitgewerkt (met medesturent VOORNAAM NAAM als mede-auteur).
% 

\section{Inleiding}%
\label{sec:inleiding}
Binnen Proximus ondersteund de Genesys-omgeving de afhandeling van klantinteracties in het callcenter en communicatie platform. Deze interacties verlopen via meerde kanalen en kunnen volgens uiteenlopende scenario's evolueren, afhankelijk van de gegevens die een klant aanlevert of opgehaald worden intern. Hierdoor ontstaat een complex proces waarin vergelijkbare gesprekken niet noodzakelijk dezelfde stappen doorlopen. Dit maakt het waardevol om de volledige conversatiestroom systematisch in kaart te brengen en te analyseren.

Hoewel Proximus uitgebreide logging verzamelt in Splunk is het vandaag niet vanzelfsprekend om deze data te vertalen naar een heldere overzicht van de volledige interacties. De beschikbare logbestanden bevatten wel veel informatie over afzondelijke uitvoeringsstappen maar zijn operationeel van aard en niet ontworpen om automatisch de volledige flowstructuur te reconstrueren. Daardoor blijft ondeuidelijk waar klanten afhaken, waar fouten optreden en welke stappen structureel inefficiënt. Een probleem dat ook wordt bevestigd in academische litartuur over IVR-analyse en process mining \autocite{Ellway2016,Antonio2024}. Dit onderzoek vertrekt daarom vanuit de vraag hoe zulke interacties beter zichtbaar en interpreteerbaar kunnen worden gemaakt.

\paragraph{Centrale onderzoeksvraag}
Hoe kan een reproduceerbare methode worden ontwikkeld om alle mogelijke paden binnen de lokale Gensys-flows van Proximus te visualiseren en te analyseren zodat drop-offpunten, foutpadden en inefficiënte stappen kunnen worden geïdentificeerd?
\paragraph{Onderzoeksdoelstelling}
Het doel van dit onderzoek is het ontwikkelen van een proof-of-concept dat Splunk-logs omzet in gereconstrueerde klantpadden binnen de Genesys-omgeving van Proximus. Deze paden worden vervolgens geclusterd tot standaardflows zodat afwijkingen, ineffici*enties en foutpaden kunnen worden gedetecteerd. Het prototype moet inzicht bieden in de volledige flowstuctuur en ondersteuning bieden bij het verbeteren van efficiëntie en nauwkeurigheid.
\paragraph{Deelvragen probleemdomein}
\begin{itemize}
  \item Welke soorten interacties en events worden door Proximus gelogd in Splunk en hoe verhouden deze zich tot de onderliggende Genesys-flows?
  \item Welke datakwaliteitsproblemen (zoals ontbrekende events, inconsistenties of timingverschillen) bemoeilijken het reconstrueren van klantpaden?
\end{itemize}

\paragraph{Deelvragen oplossingsdomein}
\begin{itemize}
  \item Hoe kunnen Splunk-logs gestructureerd en verwerkt worden om volledige klantpaden automatisch te reconstrueren?
  \item Welke analysetechnieken of algoritmen zijn geschikt om alle mogelijke paden binnen een Genesys-flow te reconstrueren?
  \item Hoe kunnen afwijkende paden, foutpaden en inefficiënte stappen automatisch worden gedetecteerd op basis van clusters?
  \item Welke visualisatietechnieken zijn het meest geschikt om zowel standaardflows als afwijkende paden inzichtelijk te maken?
\end{itemize}

%---------- Stand van zaken ---------------------------------------------------


\section{Literatuurstudie}%
\label{sec:literatuurstudie}

% Hier beschrijf je de \emph{state-of-the-art} rondom je gekozen onderzoeksdomein, d.w.z.\ een inleidende, doorlopende tekst over het onderzoeksdomein van je bachelorproef. Je steunt daarbij heel sterk op de professionele \emph{vakliteratuur}, en niet zozeer op populariserende teksten voor een breed publiek. Wat is de huidige stand van zaken in dit domein, en wat zijn nog eventuele open vragen (die misschien de aanleiding waren tot je onderzoeksvraag!)?

% Je mag de titel van deze sectie ook aanpassen (literatuurstudie, stand van zaken, enz.). Zijn er al gelijkaardige onderzoeken gevoerd? Wat concluderen ze? Wat is het verschil met jouw onderzoek?

% Verwijs bij elke introductie van een term of bewering over het domein naar de vakliteratuur, bijvoorbeeld~\autocite{Hykes2013}! Denk zeker goed na welke werken je refereert en waarom.

% Draag zorg voor correcte literatuurverwijzingen! Een bronvermelding hoort thuis \emph{binnen} de zin waar je je op die bron baseert, dus niet er buiten! Maak meteen een verwijzing als je gebruik maakt van een bron. Doe dit dus \emph{niet} aan het einde van een lange paragraaf. Baseer nooit teveel aansluitende tekst op eenzelfde bron.

% Als je informatie over bronnen verzamelt in JabRef, zorg er dan voor dat alle nodige info aanwezig is om de bron terug te vinden (zoals uitvoerig besproken in de lessen Research Methods).

% % Voor literatuurverwijzingen zijn er twee belangrijke commando's:
% % \autocite{KEY} => (Auteur, jaartal) Gebruik dit als de naam van de auteur
% %   geen onderdeel is van de zin.
% % \textcite{KEY} => Auteur (jaartal)  Gebruik dit als de auteursnaam wel een
% %   functie heeft in de zin (bv. ``Uit onderzoek door Doll Hill (1954) bleek
% %   ...'')

% Je mag deze sectie nog verder onderverdelen in subsecties als dit de structuur van de tekst kan verduidelijken.

%---------- Methodologie ------------------------------------------------------
\section{Methodologie}%
\label{sec:methodologie}
Dit onderzoek volgt een experimentele en data-gedreven aanpak. Het doel is om Splunk-logs om te zetten in werkelijke klantpaden, deze paden te clusteren tot standaardflows en vervolgens afwijkingen en inefficiënties te analyseren. 
\paragraph{Literatuurstudie}
Om de vragen in het probleemdomein te beantwoorden wordt er een gerichte literatuurstudie uitgevoerd naar:
\begin{itemize}
  \item IVR-analyse en drop-offdetectie
  \item process mining en padreconstructie
  \item log-analyse in contactcenteromgevingen
  \item clusteringtechnieken voor sequentiële data
  \item grafmodellen voor flowvisualisatie
\end{itemize}
Deze stude dient om bestaande methoden en relevante concepten te identificeren. De literatuur vormt de theoretische basis voor de ontwerpkeuzes in het proof-of-concept.
\paragraph{Analyse en voorbereiding van Splunk-logs}
De Splunk-logs van Proximus worden geanalyseerd om te bepalen:
\begin{itemize}
  \item welke events relevant zijn voor flowstappen
  \item hoe events kunnen worden gegroepeerd tot klanttransacites
  \item welke datakwaliteitsproblemen aanwezig zijn
\end{itemize}




\paragraph{Ontwikkeling van een flowreconstructiemethode (Proof-of-Concept)}
In deze fase wordt een proof-of-concept ontwikkeld dat:
\begin{itemize}
  \item Werkelijke klantpadden reconstrueert uit Spluk-logs
  \item paden clustert tot standaardflows
  \item afwijkende paden, foutpaden en inefficiënte stappen visualiseerd
\end{itemize}
Onderzoekstechnieken:
\begin{itemize}
  \item algoritmen-ontwerp voor padreconstructie
  \item grafmodellering
  \item clustering van sequenties
\end{itemize}

\paragraph{Validatie}
De ontwikkelde methode wordt geëvalueerd op basis van correctheid, compleetheid en bruikbaarheid.

\paragraph{Tools en technologieën}
Het onderzoek maakt gebruik van:
\begin{itemize}
  \item Cisco Splunk voor logverwerking
  \item Python
  \item Jupyter Notebooks voor experimentatie
  \item Git voor versiebeheer
\end{itemize}

\paragraph{Tijdsplanning en deliverables (13 weken)}
\begin{itemize}
  \item \textbf{Week 1--2:} Literatuurstudie
  \item \textbf{Week 3:} Verzameling en eerste verkenning van Splunk-logs.
  \item \textbf{Week 4--5:} Datapreprocessing en structurering van logs tot klanttransacties.
  \item \textbf{Week 6--7:} Ontwikkeling van algoritmen voor padreconstructie.
  \item \textbf{Week 8:} Clustering van gereconstrueerde klantpaden tot standaardflows.
  \item \textbf{Week 9:} Visualisatie van Standaardflows en afwijkende paden.
  \item \textbf{Week 10:} Validatie met Proximus medewerker.
  \item \textbf{Week 11-12:} Uitwerking van het volledige proof-of-concept.
  \item \textbf{Week 13:} Eindpresentatie en afronding van de bachelorproef.
\end{itemize}
%---------- Verwachte resultaten ----------------------------------------------
\section{Verwacht resultaat, conclusie}%
\label{sec:verwachte_resultaten}

